\section{Conclusions}
\subsection{CTR improvement}
Waveform-based correction substantially improved leading-edge 
timing on the first-stage energy signals of both boards. 
The test CTR was reduced by more than $20\%$ 
even with the short Onishi-inspired window, with reductions 
reaching approximately $25\%$ when the wider window was used. 
The strong performance of the
Onishi-inspired CNN in the short-window energy comparisons is
consistent with previous waveform-based LED correction
studies~\cite{onishi2022,loignon2025}.

\subsection{Shared detector-level models}
The detector-shared formulation retained competitive timing
performance without using separate detector-specific scoring
functions. Across the eight board, mode, and window comparisons, 
the best detector-shared estimator achieved a test CTR within 
$2\%$ of the overall best in six cases, with a maximum relative 
difference of approximately $3.5\%$. 
This shows that most of the achievable timing correction can be 
retained using a shared detector-level model, without requiring 
separate parameters for each detector.

\subsection{Model comparison}
CNN and MLP approaches did not consistently outperform linear 
regression, despite the successful application of CNNs in previous 
waveform-based timing studies~\cite{onishi2022,loignon2025}. 
In all four wide-window comparisons, the lowest test CTR 
was achieved either by linear ridge regression or by MiniRocket 
features followed by ridge regression. 
This suggests that much of the waveform-dependent timing 
correction can be captured by linear regression or linear 
combinations of extracted features, without requiring 
end-to-end nonlinear neural networks. The marked FBK
energy-wide shared-CNN failure also highlights a practical advantage
of the ridge-based estimators: for fixed input features, ridge
regression has a convex objective and is not subject to the
initialization-dependent optimization of the neural models. Together
with their competitive timing performance, this makes the linear
approaches particularly attractive for this correction task.

\subsection{Window length and XAI}
Extending the waveform window consistently improved
 the best timing resolution on both boards, 
 with a particularly important effect on the timing signals, 
 where short windows provided limited correction. 
 Temporal occlusion analysis showed that energy-channel 
 models relied primarily on the region around the LED crossing,  
 while timing-channel models were also sensitive to later waveform structures. 
 This was especially evident for the FBK board, where a 
 pronounced secondary sensitivity peak appeared around $20\,\mathrm{ns}$. 
 These findings support the usefulness of waveform information 
 beyond the rising edge.

 \subsection{Final considerations}
 Despite ML correction, the second-stage timing signals achieved 
 approximately twice the timing resolution of the first-stage energy signals, 
 highlighting that waveform-based correction cannot fully compensate for differences 
 in analog readout quality. 
 Nevertheless, ML consistently improved LED timing, with both boards reaching similar 
 corrected CTRs of approximately $53$–$54\,\mathrm{ps}$ despite different amplification 
 electronics and LED thresholds. 
 This convergence suggests that the remaining timing uncertainty may be dominated by 
 physical effects common to the detectors that cannot be recovered from the available 
 waveform information, although this hypothesis requires further investigation.
